\documentclass{situationist-preprint}

\title{The Perfect Forgery}
\msid{AES-04-a}
\msfamily{AES-04}
\msdiscipline{Aesthetics}
\mstype{Principal manuscript}
\mskeywords{forgery; originality; aesthetic value; artistic value; achievement; authenticity; art market}

\begin{document}
\maketitle

\begin{abstract}
A perfect forgery is indistinguishable from the original by any inspection. If the value of a painting lies in how it looks, the forgery is as valuable as the original, and the large difference in their prices is a symptom of snobbery. We argue that the forgery is not as valuable as the original, while granting the premise about looks. We distinguish aesthetic value, which supervenes on the perceptible features of an object in its category, from artistic value, which includes the value of the achievement the work embodies and depends on the work's history. We show that a perfect forgery and its original have the same aesthetic value but differ in artistic value, because the forgery does not embody the achievement that the original does. We then argue that the difference in market price is best explained by the difference in artistic value: the premium follows attributions of achievement, and it does not follow either appearance or mere physical contact with the artist. The result reconciles Lessing's claim that the forgery is aesthetically as good with Dutton's claim that it misrepresents an achievement.
\end{abstract}

\section{The problem}

In 1945 Han van Meegeren, accused of having sold a Vermeer to Hermann G\"oring, proved that the painting was his own work by painting another \enquote{Vermeer} under supervision. Paintings that had been praised by leading critics as masterpieces of Vermeer's early period lost most of their value overnight, although nothing about their appearance had changed. The case poses a question that does not depend on its details: if a forgery is perfect, so that no inspection could tell it from the original, is it as valuable as the original?

\citet{Lessing1965} answered that it is, aesthetically. On his view, what the forgery lacks is originality, which belongs to the artist and to the history of art rather than to the work as an object of aesthetic attention; to value the original above the perfect copy on aesthetic grounds is to confuse aesthetic with historical or financial considerations. \citet{Goodman1968} argued that there is an aesthetic difference between a forgery and an original even when they cannot at present be told apart by looking, because the knowledge that they differ affects how we look at them and may train us to perceive differences we cannot yet perceive. \citet{Dutton1979} argued that a work of art is the product of a performance, an achievement of its maker, and that a forgery misrepresents the achievement it presents; the forgery is therefore not merely historically but artistically deficient. \citet{Walton1970} showed that the aesthetic properties of a work depend on the category in which it is correctly perceived, and that facts about a work's origin bear on which category is correct. The views of \citet{Currie1989} and \citet{DDavies2004}, on which the artwork is an action or performance rather than a physical object, develop Dutton's position into an ontology of art.

We argue for a position that keeps the strongest point on each side. Lessing is right that the perfect forgery is aesthetically identical to the original. Dutton is right that it is not as valuable. The two claims are consistent once aesthetic value is distinguished from artistic value, and the difference in price tracks the latter.

\section{Framework}

\begin{definition}[Perfect forgery]\label{d:forg}
An object $F$ is a \emph{perfect forgery} of an artwork $O$ if $F$ was made, by someone other than $O$'s maker, so as to be taken for $O$ or for another work by $O$'s maker, and $F$ is indistinguishable from $O$ by any perceptual inspection, by any observer, at any time.
\end{definition}

The last clause is stronger than anything real forgeries achieve, and is meant to be. It answers Goodman's point in advance: if a difference could at some time be detected by inspection, the forgery is not perfect in our sense, and whatever aesthetic difference that detection reveals is not a counterexample to the claims below. Definition~\ref{d:forg} covers both the copy of an existing work and the pastiche, such as van Meegeren's, that imitates an artist's manner in a new composition. In the second case, the comparison is between the forgery and the work it purports to be: a painting by Vermeer with exactly its perceptible features.

\begin{definition}[Aesthetic and artistic value]\label{d:values}
The \emph{aesthetic value} of a work is the value it has in virtue of its aesthetic properties, such as beauty, gracefulness, unity and expressiveness. Its \emph{artistic value} is the value it has as a work of art, which includes its aesthetic value together with the value of the achievement it embodies: the solution of artistic problems, originality of conception, and the skill exercised in realizing it.
\end{definition}

\begin{principle}[Aesthetic supervenience]\label{p:sup}
The aesthetic properties of a work supervene on its perceptible properties together with the category in which it is correctly perceived.
\end{principle}

\begin{principle}[Achievement]\label{p:ach}
Artistic value includes achievement value: other things equal, a work that embodies a greater achievement by its maker has greater artistic value. Achievement is relative to the maker's situation, including what had been done before and what problems the maker solved.
\end{principle}

\begin{principle}[Categories]\label{p:cat}
The categories relevant to Principle~\ref{p:sup} are perceptually distinguishable categories, such as medium, genre and style: classes of works whose membership is marked by perceptible features \citep{Walton1970}.
\end{principle}

Principle~\ref{p:sup} is the thesis that aesthetic properties depend on how a work looks or sounds, as heard or seen in the right way. Principle~\ref{p:cat} follows Walton's definition: the categories that shape aesthetic perception are those whose membership is determined by perceptible features standard, variable or contra-standard for them. Origin can decide which such category is correct for a work, but the categories themselves are perceptual. Principle~\ref{p:ach} is Dutton's thesis that we appreciate works of art, in part, as achievements \citep{Dutton1979}. It explains why we admire a painter who solved a problem of representing light for the first time more than one who applied the solution later, why a performance of a difficult piece is admired for its difficulty, and why the discovery that a work was made by mechanical means alters our assessment of it.

\section{The argument}

\begin{lemma}\label{l:aesth}
A perfect forgery $F$ of $O$ has the same aesthetic properties, and hence the same aesthetic value, as $O$.
\end{lemma}

\begin{proof}
By Definition~\ref{d:forg}, $F$ and $O$ do not differ in any perceptible property. Their correct perceptual categories are the same: the categories relevant to aesthetic properties are perceptually distinguishable (Principle~\ref{p:cat}), and $F$ was made to have, and does have, exactly the perceptible features that mark $O$'s membership in its medium, genre and style. By Principle~\ref{p:sup}, $F$ and $O$ have the same aesthetic properties. Aesthetic value is the value a work has in virtue of its aesthetic properties (Definition~\ref{d:values}), so they have the same aesthetic value.
\end{proof}

\begin{lemma}\label{l:art}
A perfect forgery $F$ of $O$ embodies a lesser achievement than $O$, in the respects that make up achievement value.
\end{lemma}

\begin{proof}
The achievement embodied in $O$ includes the conception of its composition, the invention or development of the manner in which it is painted, and its execution, relative to the art of its time (Principle~\ref{p:ach}). If $F$ is a copy of $O$, its maker did not conceive the composition or invent the manner, and executed a design already given. If $F$ is a pastiche, its maker conceived a composition in a manner already invented by another, and in a situation in which that manner was no longer a solution to any open problem. In both cases the forger's achievement is one of imitation, which may be considerable, but it lacks the components of originality that $O$'s achievement includes. Moreover, $F$ presents itself as embodying $O$'s maker's achievement, which it does not.
\end{proof}

\begin{theorem}\label{t:main}
A perfect forgery is aesthetically identical to its original but has less artistic value. It is therefore not as valuable as the original.
\end{theorem}

\begin{proof}
By Lemma~\ref{l:aesth}, the two have the same aesthetic value. By Lemma~\ref{l:art} and Principle~\ref{p:ach}, the forgery has less achievement value. Artistic value includes both (Definition~\ref{d:values}), so the forgery has less artistic value, and since its value as a work of art is its artistic value, it is less valuable as a work of art.
\end{proof}

Theorem~\ref{t:main} explains the van Meegeren case without imputing error to anyone's eyes. The critics who praised the forged Vermeers were not wrong about how the paintings looked, and the paintings did not come to look worse. What changed was the achievement they could be credited with: a painting thought to show the young Vermeer working out a religious manner turned out to show a twentieth-century painter imitating a manner already known. The case is examined in detail in the companion manuscript AES-04-b.

\section{What the price tracks}

The market price of a perfect forgery, once known as such, is a small fraction of the price of the original. We now argue that this difference tracks artistic value rather than any of three alternatives.

\begin{proposition}\label{p:price}
The price difference between an original and a perfect forgery of it is explained by the difference in artistic value, and not by a difference in aesthetic value, in physical contact with the artist, or in mere scarcity.
\end{proposition}

\begin{proof}
\emph{Not aesthetic value.} By Lemma~\ref{l:aesth}, there is no difference in aesthetic value to explain it.

\emph{Not contact.} If the premium were for objects the artist had touched, it would attach equally to every object from the artist's studio, and it does not: a canvas the artist primed and left blank, or a palette, commands a fraction of the price of a finished work. Conversely, the premium attaches to achievement where contact is absent or shared. Works produced in a master's workshop are priced by the share of the master's own invention and hand in them, and a copy after another painter made by a major artist, such as Rubens's copies after Titian, is priced as a work by that artist, since it embodies that artist's achievement in interpreting the original.

\emph{Not mere scarcity.} A perfect forgery is as unique as its original: there is one of each. Scarcity explains why originals are expensive only given that what is scarce is valued, and what is valued is the work as the record of a particular achievement.

\emph{Artistic value.} The premium follows attributions of achievement in each of these cases: it rises and falls with the attribution of the work's conception and execution to the artist whose achievement it is supposed to embody, and it falls when an attribution is withdrawn although nothing perceptible has changed. This is the pattern that the difference in artistic value predicts.
\end{proof}

Empirical work on judgments of value supports the conclusion. \citet{NewmanBloom2012} found that participants value original artworks above identical duplicates, and that the premium depends both on the original's being the product of a creative act and on its contact with the artist; Proposition~\ref{p:price} does not require that contact plays no role in individual valuations, only that it does not explain the structure of the premium described above.

\section{Objections and replies}

\begin{objection}[Goodman]\label{o:good}
Knowing that two pictures differ, even if we cannot now see the difference, gives us reason to look at them differently and may train us to see a difference. The aesthetic properties of the original are therefore not exhausted by what can now be seen \citep{Goodman1968}.
\end{objection}

\begin{reply}
Goodman's argument shows that the perceptual indiscernibility of real forgeries is relative to the training of the observer, and that real forgeries are almost never perfect. It does not show that a forgery indiscernible from its original by any observer at any time differs aesthetically. Definition~\ref{d:forg} stipulates the stronger case. If no training could reveal a difference, the knowledge that the pictures differ in origin may change how we look, but not what there is to be seen.
\end{reply}

\begin{objection}[Categories]\label{o:walton}
By Walton's own argument, a work's history bears on its correct category. The original Vermeer is correctly perceived as a seventeenth-century Dutch painting; the van Meegeren as a twentieth-century pastiche. Their aesthetic properties therefore differ \citep{Walton1970}.
\end{objection}

\begin{reply}
History bears on which perceptual category is correct; it does not create categories that are not perceptual. \enquote{Twentieth-century pastiche of Vermeer} is not marked by any perceptible features that the perfect forgery has and the original lacks, since by hypothesis there are none, and the standard, variable and contra-standard features of the forgery relative to Vermeer's style are exactly those of the original. Where a forgery is detectable as pastiche, for example by an anachronism of manner, it is not perfect. The differences the objection points to are real, but they are differences in achievement and in historical position, which is what Lemma~\ref{l:art} records.
\end{reply}

\begin{objection}[Aestheticism]\label{o:formal}
The value of a work of art as art is its aesthetic value. Achievement value belongs to the biography of the artist, and to prefer the original on its account is snobbery, as Lessing held \citep{Lessing1965}.
\end{objection}

\begin{reply}
Aestheticism cannot explain features of artistic evaluation that it does not regard as confused. We rank a work that first achieved a mode of representation above later works indistinguishable in quality; we admire a performance for the difficulty it overcomes; we judge conceptual works, which may have no distinctive aesthetic properties at all, for what they achieve. On the institutional account of art defended in AES-01, a work's status as art itself depends on its relation to an artworld, so that art-relevant properties are not confined to perceptible ones. Achievement value is part of what we value in works as art, and to exclude it is to stipulate a narrower notion than the practice uses.
\end{reply}

\begin{objection}[Prices]\label{o:market}
Art prices are driven by fashion, speculation and status. It is implausible that they track any kind of value.
\end{objection}

\begin{reply}
Proposition~\ref{p:price} does not claim that prices measure artistic value, only that the \emph{difference} between the price of an original and of a perfect forgery tracks the difference in artistic value. Fashion and speculation affect both prices. The structure of the difference, its attachment to attributions of achievement and its indifference to appearance and to contact as such, is what the proposition explains, and it is what the snobbery hypothesis leaves unexplained, since snobbery about origin would attach to contact as much as to achievement.
\end{reply}

\section{What the result does not establish}

The argument does not show that forgeries are without value. A perfect forgery has all the aesthetic value of its original and may embody considerable achievement of its own as an imitation. Nor does it address the moral wrong of forgery, which consists in deception and fraud and is independent of the comparison of values. It is limited to singular arts. In arts with multiple instances, such as printmaking, cast sculpture or music, an indiscernible instance produced in the authorized way is not a forgery but another instance of the work, and the question of forgery attaches to the provenance of the instance rather than to its appearance. Finally, the argument does not show that every difference of history makes a difference to artistic value: the claim is restricted to differences in the achievement the work embodies.

\section{Conclusion}

A perfect forgery is not as valuable as its original. It looks exactly the same and has the same aesthetic value, as Lessing held; but it does not embody the original's achievement, and the value of a work as art includes the value of the achievement it embodies, as Dutton held. The difference in price between an original and its perfect forgery tracks this difference in artistic value: it follows attributions of achievement, and neither appearance nor contact with the artist.

\begin{thebibliography}{99}\small
\bibitem[Currie(1989)]{Currie1989} Currie, G. (1989). \emph{An Ontology of Art}. Macmillan.
\bibitem[Davies(2004)]{DDavies2004} Davies, D. (2004). \emph{Art as Performance}. Blackwell.
\bibitem[Dutton(1979)]{Dutton1979} Dutton, D. (1979). Artistic crimes: The problem of forgery in the arts. \emph{The British Journal of Aesthetics}, 19(4), 302--314.
\bibitem[Goodman(1968)]{Goodman1968} Goodman, N. (1968). \emph{Languages of Art: An Approach to a Theory of Symbols}. Bobbs-Merrill.
\bibitem[Lessing(1965)]{Lessing1965} Lessing, A. (1965). What is wrong with a forgery? \emph{The Journal of Aesthetics and Art Criticism}, 23(4), 461--471.
\bibitem[Newman and Bloom(2012)]{NewmanBloom2012} Newman, G.~E. and Bloom, P. (2012). Art and authenticity: The importance of originals in judgments of value. \emph{Journal of Experimental Psychology: General}, 141(3), 558--569.
\bibitem[Walton(1970)]{Walton1970} Walton, K.~L. (1970). Categories of art. \emph{The Philosophical Review}, 79(3), 334--367.
\end{thebibliography}

\end{document}
